% Part III: connecting and building databases. Sources: docs/databases.md ("Choose a way in",
% "What every database must say", "Connect an OMOP database", "Connect any SQL database through
% views" with its "Canonical views reference", "Build a new database for Medorama", "Files:
% recordings, ECGs and images", "Check the result"); docs/design/sources.md;
% docs/administration.md, "Connected databases". Checked against medorama/src/sources/mod.rs
% (connections, settings, views, checks, filters), sources/omop.rs (the OMOP views),
% sources/link.rs (the registry), categories.rs (of_entry), main.rs (connections commands),
% server.rs (connection_add) and web/index.html (the form).
\chapter{Connecting and building databases}
\label{chap:connecting}

A \emph{connected database} is a database that Medorama reads in place, read-only, at each
request, beside the catalog: an OMOP CDM instance, a research warehouse, a registry or a study's
database. Nothing is copied from it but its persons' identifiers, once they are linked to the
catalog's patients.

This chapter is written for data managers and database engineers. It explains which way into
Medorama to choose, what every database must say, how to connect an OMOP database and any other
SQL database, how to build a new database that Medorama reads, and how a connected database's
persons become the catalog's patients. Importing tables is the subject of \cref{chap:importing}.

\section{Three ways in}
\label{sec:connecting-ways}

\begin{table}[htbp]
\caption{The three ways a database reaches Medorama.}
\label{tab:connecting-ways}
\begin{tabularx}{\linewidth}{@{}P{2.4cm}LLL@{}}
\toprule
 & {\tabhead Import tables} & {\tabhead Connect an OMOP database} & {\tabhead Connect any other
  database} \\
\midrule
What & CSV, Parquet, Excel or OpenDocument files, CLIF folders, known exports, and the tables
  that \cmd{muse-py index} and \cmd{dicom-py index} write & An OMOP CDM 5.3 or 5.4 instance & A
  research warehouse, a registry or a study's database, in any layout \\
Where the data live & Copied into the catalog: each table is a database of the Databases tab,
  kept as it came & In place: read at each request, never copied (only persons' identifiers, once
  linked) & In place, as for OMOP \\
Speed and freshness & Fastest searches; as current as the last import & Each search queries the
  database; always current & The same as OMOP \\
Engines & None needed & PostgreSQL (on-site, Amazon RDS and Aurora, Azure), SQLite files & The
  same as OMOP \\
What you write & Nothing, if the column names are recognized; otherwise a mapping & The
  connection's settings & Three SQL views (\code{persons}, \code{visits}, \code{entries}) in a
  TOML file \\
Files (recordings, ECGs, images) & Yes, through a \code{file} column & No & No \\
Removing & Removes everything it added & Disconnecting removes its links and what it added; the
  database is untouched & The same as OMOP \\
\bottomrule
\end{tabularx}
\end{table}

To choose, answer these questions in order:
\begin{enumerate}
\item \textbf{Are they files} (recordings, ECGs, DICOM)? Then list them in a table with their
  patients and import it. A connected database cannot bring files, and a file never enters
  without its patient.
\item \textbf{Is it an OMOP CDM database?} Connect it as \code{omop}, rather than exporting it.
\item \textbf{Must the data stay in their database}, because they change daily or are too large
  to copy? Connect it through views.
\item \textbf{Is it a one-off extract, or a set of spreadsheets?} Import it.
\end{enumerate}
The three ways mix in one catalog. A person of a connected database is linked to the same patient
as an imported row with the same patient ID or MRN (\cref{sec:connecting-registry}), so one
patient's page can show imported laboratory results, an OMOP instance's diagnoses and a recording
side by side.

\screenshot{sources.pdf}{Sources: the catalog and the connected databases answer the same questions; their persons are linked to one registry of patients, and every answer leaves through the privacy gate.}{fig:connecting-sources}

\section{What every database must say}
\label{sec:connecting-must-say}

Whatever the way in, the data must answer five questions for Medorama to place them.
\begin{enumerate}
\item \textbf{Who.} Every row names its patient by an identifier that the institution uses
  consistently: a \code{patient\_id} (a study's or a system's ID), or an \code{mrn} (the medical
  record number), with \code{mrn\_authority} when several systems issue MRNs. Names, date of
  birth and sex describe a patient; they never identify one on their own.
\item \textbf{When.} Every result, note, drug or file has a time: with an offset
  (\code{2024-03-02T07:10:00-05:00}, \code{Z}), as local time in the catalog's or the
  connection's time zone (\code{2024-03-02 07:10}), or in Unix seconds.
\item \textbf{Where in the stay} (optional, recommended). The encounter or visit that the row
  belongs to and, for imports, the bed stays (\code{care\_unit}, \code{bed}, \code{bed\_start},
  \code{bed\_end}) that link bedside-monitor recordings to their encounters.
\item \textbf{What.} A name (\code{Lactate}, \code{QTc}, \code{Vancomycin}) and, when there is
  one, a value with its unit, a normal range and a flag, a code with its vocabulary
  (\code{2524-7}, \code{LOINC}), or a text.
\item \textbf{Which category.} The column of the Patient Journey that it belongs to, given by its
  key (\cref{tab:connecting-keys}). The keys include \code{eeg} and \code{sleep}, the
  categories of EEG recordings and sleep studies.
\end{enumerate}

\begin{table}[htbp]
\caption{The category keys.}
\label{tab:connecting-keys}
\begin{tabularx}{\linewidth}{@{}P{1.3cm}LP{1.3cm}LP{1.3cm}LP{1.3cm}L@{}}
\toprule
{\tabhead Key} & {\tabhead Category} & {\tabhead Key} & {\tabhead Category} & {\tabhead Key} &
  {\tabhead Category} & {\tabhead Key} & {\tabhead Category} \\
\midrule
\code{icu} & ICU & \code{vit} & Vitals & \code{lab} & Labs & \code{micro} & Microbiology \\
\code{path} & Pathology & \code{ecg} & ECG & \code{holter} & Holter & \code{eeg} & EEG \\
\code{sleep} & Sleep & \code{echo} & Echo & \code{cath} & Cath lab & \code{img} & Imaging \\
\code{proc} & Procedures & \code{med} & Medications & \code{blood} & Transfusions &
  \code{support} & Life support \\
\code{dx} & Conditions & \code{obs} & Assessments & \code{note} & Notes & \code{allergy} &
  Allergies \\
\bottomrule
\end{tabularx}
\end{table}

A connected database's entry may also give a category by its label or one of its words, as an
imported row may (\cref{chap:importing}), but storing the keys themselves, in lower case, is the
safest choice. A value that names no category is read as \code{lab}. An entry of \code{lab} is
read as \code{vit} when it is a vital sign, recognized by its name (\code{heart rate},
\code{SpO2}, \code{MAP} and so on) or by its LOINC code (\code{8867-4}, \code{8310-5},
\code{8480-6}, \code{8462-4}, \code{8478-0}, \code{9279-1}, \code{59408-5}, \code{2708-6},
\code{29463-7}, \code{8302-2}, \code{39156-5}, \code{8331-1}, \code{8332-9}, \code{8333-7}).

\begin{phibox}
Identifiers are protected data. Medorama keeps names, MRNs, account numbers and dates of birth in
its identifier tables. Users without the PHI permission, and everyone in Privacy Mode, see
placeholders and pseudonyms instead (\cref{chap:protecting}). Export or expose only the
identifiers needed for linking: a study that needs no names or dates of birth should not provide
them, since Medorama links by patient IDs and MRNs alone.
\end{phibox}

\section{A connection}
\label{sec:connecting-connection}

A connection is a row of the catalog's database: a name, a kind, the database's URL without its
password, the password's file, and settings (\cref{tab:connecting-settings}). Only
administrators connect and disconnect databases. Connecting (\code{connect}), linking and
syncing (\code{link}: the persons linked, held and invalid) and disconnecting
(\code{disconnect}) are written to the audit log.

\begin{itemize}
\item \textbf{Kinds.} The kind \code{omop} is an OMOP CDM 5.3 or 5.4 database, whose views are
  built in (\cref{sec:connecting-omop}). The kind \code{sql} is any layout, read through views
  that an administrator writes (\cref{sec:connecting-sql}).
\item \textbf{Databases.} Medorama reads PostgreSQL (on-site, Amazon RDS and Aurora, Azure
  Database for PostgreSQL, anything that speaks its protocol), given by a \code{postgres://} or
  \code{postgresql://} URL, and SQLite files, given by a \code{sqlite:} URL
  (\code{sqlite:/data/county/omop.sqlite}). Other engines (SQL Server, Oracle, Snowflake,
  BigQuery) are not read directly yet; a PostgreSQL foreign-data wrapper reaches them in the
  meantime.
\item \textbf{Read only, enforced.} PostgreSQL sessions are opened with
  \code{default\_transaction\_read\_only = on}, and SQLite files are opened read-only, for
  queries only. Medorama therefore cannot write there, even with an account that could. Give it
  an account that can only read as well.
\item \textbf{The name} is written in criteria before a dot (\code{omop.lactate}). It is made of
  letters, digits, \code{-} and \code{\_}, starts with a letter, has at most 40 characters, and
  is not the name of another connection or of an imported database.
\item \textbf{Secrets.} The password is read from a file (\opt{-{}-password-file}) or from the
  environment variable \code{MEDORAMA\_SOURCE\_PASSWORD\_<NAME>}, which takes precedence. In the
  variable's name, the connection's name is written in capitals, with every character other than
  letters and digits written as \code{\_}: \code{omop-cdm} reads
  \code{MEDORAMA\_SOURCE\_PASSWORD\_OMOP\_CDM}. A password given in the page is kept in
  \file{<home>/secrets/<name>.password}, readable by its owner only, and deleted when the
  database is disconnected. A password is never stored in the database, in a settings file, in a
  URL or in the audit log, and it is never sent back. A URL that holds a password is refused.
\item \textbf{Checked when added.} Each view must answer with its canonical columns, and the
  version query (if any) with a value; otherwise the connection is refused.
\end{itemize}

\begin{xltabular}{\linewidth}{@{}P{3.6cm}P{2.6cm}L@{}}
\caption{The settings of a connection (\cmd{medorama connections add}).}
\label{tab:connecting-settings}\\
\toprule
\tabhead Option & \tabhead Default & \tabhead Meaning \\
\midrule
\endfirsthead
\toprule
\tabhead Option & \tabhead Default & \tabhead Meaning \\
\midrule
\endhead
\bottomrule
\endfoot
\code{NAME} & (Required) & The connection's name, as criteria write it \\
\opt{-{}-kind} & (Required) & \code{omop} or \code{sql} \\
\opt{-{}-url} & (Required) & \code{postgres://user@host/db?sslmode=require} (without a password),
  or \code{sqlite:/path/to/file} \\
\opt{-{}-password-file} & (None) & The file that holds the password; otherwise the variable
  \code{MEDORAMA\_SOURCE\_PASSWORD\_<NAME>} \\
\opt{-{}-ca-file} & (None) & Certificate authorities (PEM) for the server's certificate, besides
  the system's \\
\opt{-{}-schema} & The default schema & The schema of the database's tables (OMOP's CDM): a
  plain name (letters, digits, \code{\_}) \\
\opt{-{}-vocabulary-schema} & The CDM's & The schema of OMOP's vocabulary (\code{concept}), when
  it is separate: a plain name \\
\opt{-{}-timezone} & The catalog's & The time zone of the database's local times, an IANA name
  such as \code{Europe/Paris} (an unknown one is refused). Without this option, the catalog's
  time zone is written into the connection when it is added (by the command or the page) and
  stays fixed from then on: changing the catalog's \code{timezone} later does not change it. On
  PostgreSQL, a \code{timestamp} column (without a time zone) is read in it; a
  \code{timestamptz} is an instant \\
\opt{-{}-identifier} & \code{mrn} & What OMOP's \code{person\_source\_value} holds: \code{mrn},
  \code{patient\_id} or \code{none} \\
\opt{-{}-mrn-authority} & (None) & The authority that issued the database's MRNs, for persons
  whose view gives none \\
\opt{-{}-timeout-s} & 60 & Each query's time limit, in seconds (PostgreSQL only; SQLite files are
  local and have none) \\
\opt{-{}-max-rows} & 1{,}000{,}000 & The number of rows a criterion may read, and of persons a
  sync reads \\
\opt{-{}-views} & (None) & For the kind \code{sql}, a TOML file with the views \\
\end{xltabular}

\paragraph{In the page,} an administrator uses \ui{Import} $\rightarrow$ \ui{Databases}, in the
card \ui{Add a database}, and then chooses \ui{Connect a database file} (an SQLite file on the
server, inside \code{data\_roots}, never the catalog's own database; \cmd{connections add}
checks the same) or \ui{Connect a database server} (a PostgreSQL server). A name already used is
refused before anything is written (a password given in the page included). The form asks for:
\begin{itemize}
\item the \ui{Name};
\item the \ui{Layout}: \ui{OMOP CDM 5.3 or 5.4}, or \ui{Any layout, through SQL views}, which
  opens \ui{SQL views (any layout)} with the boxes \ui{persons}, \ui{visits}, \ui{entries} and
  \ui{version};
\item the \ui{File on this server}, or the \ui{Server} and its \ui{Password};
\item the \ui{Schema} and the \ui{Vocabulary schema};
\item \ui{Its patients are named by} (\ui{MRN}, \ui{patient ID} or \ui{nothing the catalog
  knows}), the \ui{MRN authority}, and the \ui{Time zone of its times}.
\end{itemize}
\ui{Connect} checks the database and connects it. The card \ui{Connected databases} lists the
connections, each with \ui{Check}, \ui{Link all} and \ui{Disconnect\dots}.

The same can be done on the command line. The following commands connect an OMOP database on a
PostgreSQL server and another in an SQLite file, then list, check, link and remove connections:
\begin{lstlisting}[style=shell]
medorama connections add omop --kind omop \
    --url "postgres://reader@omop.example.org/cdm?sslmode=require" \
    --password-file /etc/medorama/omop-password \
    --schema cdm --vocabulary-schema vocab --mrn-authority HOSP --timezone America/New_York
medorama connections add county --kind omop --url sqlite:/data/county/omop.sqlite --identifier mrn
medorama connections list            # name, kind, URL, engine, persons linked
medorama connections check omop      # each view answers; the data's version
medorama connections sync omop       # link every person now (else as searches find them)
medorama connections remove omop     # its links and what it added go; the database is untouched
\end{lstlisting}

\section{Connecting an OMOP database}
\label{sec:connecting-omop}

Medorama has the views of the OMOP CDM 5.3 and 5.4 built in. They read the tables \code{person},
\code{visit\_occurrence}, \code{visit\_detail}, \code{care\_site}, \code{measurement},
\code{observation}, \code{condition\_occurrence}, \code{drug\_exposure},
\code{procedure\_occurrence}, \code{note}, \code{provider}, \code{concept} (in the vocabulary's
schema) and \code{cdm\_source} (the version). Every one of these tables must exist, even if it is
empty: a view that names a missing table fails its check, and the connection is refused when it
is added. \code{cdm\_source} is therefore required in effect, since the version is read from it.

\begin{xltabular}{\linewidth}{@{}P{3.2cm}L@{}}
\caption{How each OMOP table is read.}
\label{tab:connecting-omop}\\
\toprule
\tabhead OMOP table & \tabhead Read as \\
\midrule
\endfirsthead
\toprule
\tabhead OMOP table & \tabhead Read as \\
\midrule
\endhead
\bottomrule
\endfoot
\code{person} & A person. \code{person\_id} is the key, and \code{person\_source\_value} is the
  MRN or the patient ID, as the setting \code{identifier} says (\code{mrn}, the default;
  \code{patient\_id}; or \code{none}, neither). The date of birth comes from
  \code{birth\_datetime}, otherwise from \code{year\_of\_birth}, \code{month\_of\_birth} and
  \code{day\_of\_birth} when all three are given. The sex comes from \code{gender\_concept\_id}
  (8532 is \code{F}, 8507 is \code{M}, anything else none). No names are read. \\
\code{visit\_occurrence} & A visit, with its start and end datetimes (otherwise dates) and its
  unit: the name of its care site (\code{care\_site}), otherwise of its visit concept. It is an
  intensive-care stay when a \code{visit\_detail} of the visit has the concept 32037 (intensive
  care), or when its unit is named as an intensive-care unit (\code{ICU}, \code{MICU},
  \code{CCU} and so on) \\
\code{measurement} & \code{lab} (\code{vit} for vital signs), with the concept's name (otherwise
  \code{measurement\_source\_value}), code and vocabulary; \code{value\_as\_number}, otherwise
  the name of \code{value\_as\_concept\_id} or \code{value\_source\_value} as a text result; the
  unit concept's code (UCUM), otherwise \code{unit\_source\_value}; and \code{range\_low} and
  \code{range\_high} \\
\code{observation} & \code{obs}, with the concept's name, code and vocabulary;
  \code{value\_as\_number}, otherwise \code{value\_as\_string} or the value concept's name; and
  the unit \\
\code{condition\_occurrence} & \code{dx}, with the concept's name; the code and vocabulary of
  the source concept (\code{condition\_source\_concept\_id}: ICD-10-CM), otherwise
  \code{condition\_source\_value}; and the start and end \\
\code{drug\_exposure} & \code{med}, with the concept's name, code and vocabulary, the start and
  end, and \code{dose\_unit\_source\_value} \\
\code{procedure\_occurrence} & \code{proc}, with the concept's name, code and vocabulary, and its
  time \\
\code{note} & \code{note}, with \code{note\_title} (otherwise the note type's name, otherwise
  \code{Note}), the author (\code{provider\_name} of its \code{provider}) and \code{note\_text} \\
\code{cdm\_source} & The version: the latest \code{source\_release\_date} and
  \code{cdm\_release\_date} \\
\end{xltabular}

Times are read as the database's local times, in the connection's \code{timezone}. On SQLite,
where OMOP's tools may store dates as Unix seconds, those are turned into text first. With
\code{identifier = none}, no person is linked to an imported patient. Each person then becomes a
patient of its own, with the patient ID \code{<connection>:<person key>}
(\cref{sec:connecting-registry}).

\section{Connecting any other database}
\label{sec:connecting-sql}

Any database that PostgreSQL or SQLite can query is connected through three \emph{canonical
views}: SQL queries over its own tables that give Medorama's canonical columns. The views are the
contract between Medorama and the database. Only these columns are read, so a view cannot pass a
field that the privacy gate does not know, and the database's other columns stay where they are.
The views are written in the database's own SQL dialect, and run as written: Medorama wraps each
one as a subquery (\code{SELECT \dots{} FROM (your view) v WHERE person\_key IN (\dots)}) and
translates nothing in it (PostgreSQL's \code{?} operators, \code{E''} strings, dollar quotes and
comments are fine). A view must be a single \code{SELECT} statement without a trailing
semicolon.

\begin{table}[htbp]
\caption{The canonical views.}
\label{tab:connecting-views}
\begin{tabularx}{\linewidth}{@{}P{1.6cm}LP{1.5cm}@{}}
\toprule
{\tabhead View} & {\tabhead Columns, in this order} & {\tabhead Required} \\
\midrule
\code{persons} & \code{person\_key}, \code{patient\_id}, \code{mrn}, \code{mrn\_authority},
  \code{last\_name}, \code{first\_name}, \code{birth\_date}, \code{sex} & Yes \\
\code{visits} & \code{visit\_key}, \code{person\_key}, \code{start\_time}, \code{end\_time},
  \code{unit}, \code{icu} & No \\
\code{entries} & \code{entry\_key}, \code{person\_key}, \code{visit\_key}, \code{category},
  \code{time}, \code{end\_time}, \code{name}, \code{code}, \code{vocabulary}, \code{value\_num},
  \code{value\_text}, \code{unit}, \code{ref\_low}, \code{ref\_high}, \code{flag},
  \code{author}, \code{text} & Yes \\
\code{version} & One value that changes when the data do (a load time, a count) & No \\
\bottomrule
\end{tabularx}
\end{table}

Each column has a \emph{privacy class}, which decides what the gate does with it
(\cref{tab:connecting-columns}).

\begin{xltabular}{\linewidth}{@{}P{3.4cm}LP{3.0cm}@{}}
\caption{The canonical columns, what they hold, and their privacy class.}
\label{tab:connecting-columns}\\
\toprule
\tabhead Column & \tabhead Holds & \tabhead Class, and in Privacy Mode \\
\midrule
\endfirsthead
\toprule
\tabhead Column & \tabhead Holds & \tabhead Class, and in Privacy Mode \\
\midrule
\endhead
\bottomrule
\endfoot
\code{person\_key}, \code{visit\_key}, \code{entry\_key} & The database's own keys, of any type
  (compared as text) & Key: never sent. Persons leave as pseudonyms, entries as tokens \\
\code{patient\_id}, \code{mrn}, \code{mrn\_authority} & What links a person to the catalog's
  patients, by the rules of an imported table. Give at least one of \code{patient\_id} and
  \code{mrn} & Identifier (\code{patient\_id}, \code{mrn}): kept in the catalog's identifier
  tables, and sent as placeholders \\
\code{last\_name}, \code{first\_name}, \code{birth\_date} (\code{YYYY-MM-DD}) & The person's
  description: empty strings when the database has none or should not give them & Identifier:
  placeholders \\
\code{sex} & \code{F}, \code{M} and so on & Clinical: as is \\
\code{start\_time}, \code{end\_time}, \code{time} & Text in a form that imports read
  (\code{2024-03-02 07:10}, ISO 8601 with an offset), in the connection's time zone, or Unix
  seconds & Date: moved back by the patient's date offset \\
\code{unit} (visits) & The ward or unit of the stay & Clinical \\
\code{icu} & 1 for an intensive-care stay, otherwise 0 (a PostgreSQL boolean works too: cohort
  searches compare \code{CAST(icu AS INT) = 1} in the database) & Clinical \\
\code{category} & A category key (\cref{tab:connecting-keys}), or a word that the categories
  recognize. An unrecognized one becomes \code{lab}, or \code{vit} when the entry's name or LOINC
  code is a vital sign's & Clinical \\
\code{name}, \code{code}, \code{vocabulary} & What the entry is: \code{Lactate}, \code{2524-7},
  \code{LOINC} & Clinical \\
\code{value\_num}, \code{value\_text}, \code{unit}, \code{ref\_low}, \code{ref\_high},
  \code{flag} & A number or a text result, its unit, its normal range and its flag (\code{H},
  \code{L}, \code{HH}, \code{LL}, \code{A} and so on) & Clinical \\
\code{author} & Who wrote a note & Staff: \code{[Staff Name]} \\
\code{text} & A note's text, which is free text & Free text: shown with the patient's own names
  and date of birth replaced, and marked (the default), or withheld when an administrator
  withholds the text of notes \\
\end{xltabular}

The views are SQL in the database's own dialect, in a TOML file given with \opt{-{}-views} (in the
page, in the four boxes of \ui{SQL views (any layout)}). Each view's columns are named by its
first query: in a \code{UNION ALL}, the later parts need no names. The following file presents a
warehouse's patients, admissions, results and notes:
\begin{lstlisting}[language=TOML]
# warehouse-views.toml
persons = '''
  SELECT p.id AS person_key, p.study_id AS patient_id, p.mrn, 'HOSP' AS mrn_authority,
         '' AS last_name, '' AS first_name, CAST(p.dob AS TEXT) AS birth_date, p.sex
  FROM research.patients p'''
visits = '''
  SELECT a.id AS visit_key, a.patient_id AS person_key,
         CAST(a.admitted AS TEXT) AS start_time, CAST(a.discharged AS TEXT) AS end_time,
         a.ward AS unit, CASE WHEN a.icu THEN 1 ELSE 0 END AS icu
  FROM research.admissions a'''
entries = '''
  SELECT 'l' || r.id AS entry_key, r.patient_id AS person_key, r.admission_id AS visit_key,
         'lab' AS category, CAST(r.collected AS TEXT) AS time, NULL AS end_time,
         r.test AS name, r.loinc AS code, 'LOINC' AS vocabulary, r.value AS value_num,
         '' AS value_text, r.unit, r.low AS ref_low, r.high AS ref_high, r.flag,
         '' AS author, '' AS text
  FROM research.results r
  UNION ALL
  SELECT 'n' || n.id, n.patient_id, n.admission_id, 'note', CAST(n.written AS TEXT), NULL,
         n.kind, '', '', NULL, '', '', NULL, NULL, '', n.author, n.body
  FROM research.notes n'''
version = "SELECT MAX(loaded_at) FROM research.loads"
\end{lstlisting}
The following commands connect the warehouse with these views and check it:
\begin{lstlisting}[style=shell]
medorama connections add warehouse --kind sql --url "postgres://reader@dw.example.org/research" \
    --password-file /etc/medorama/dw-password --views warehouse-views.toml
medorama connections check warehouse
\end{lstlisting}

\subsection{Performance}
\label{sec:connecting-performance}

Medorama asks the views only for what it needs:
\begin{itemize}
\item A patient's page reads the persons, visits and entries that it shows, with
  \code{WHERE person\_key IN (\dots)}. What was read is kept for 30 seconds, because one page
  reads it for its summary, journey and report.
\item A cohort's criterion runs one query over \code{entries}, with the filters that the database
  can apply itself: the category, the name (exact, or containing a word), a code's prefix, a
  text, and whether there is a number. Medorama itself tests values in other units,
  \code{latest}, dates and stays.
\item The names of what a database holds are read once, for criteria written as text, and kept
  for ten minutes.
\end{itemize}
Index the underlying tables on the person, the time, and the name or code by which entries are
searched. A criterion that would read more than \code{max\_rows} rows is refused, with a message
that asks to narrow it (a date range, a database). Connections are pooled per server, and a
connection that the server closed is replaced.

\section{Building a new database for Medorama}
\label{sec:connecting-building}

When a study builds its database from scratch, two \emph{lightweight databases} work without any
server: a folder of CSV files, which you import, and a single SQLite file, which you connect.
Choose by how the data will be used (\cref{tab:connecting-lightweight}).

\begin{table}[htbp]
\caption{Two lightweight databases for a study.}
\label{tab:connecting-lightweight}
\begin{tabularx}{\linewidth}{@{}P{3.6cm}LL@{}}
\toprule
 & {\tabhead A folder of CSV files} & {\tabhead An SQLite file with the canonical tables} \\
\midrule
Brings files (recordings, ECGs, images) & Yes & No (list the files in a CSV table beside it) \\
Updated by & Importing again (remove the earlier import first) & Writing to the file; Medorama
  reads it live \\
Written with & Any spreadsheet or script & Any SQLite tool or library \\
Needs & Nothing & Nothing (PostgreSQL works the same way, as a schema) \\
\bottomrule
\end{tabularx}
\end{table}

\paragraph{A folder of CSV files.} Write one CSV file per kind of table, with the column names of
\cref{tab:importing-columns}, and import them in this order:
\begin{lstlisting}[style=text]
study/
  patients.csv       patient_id, sex, birth_date (only if needed)
  encounters.csv     encounter_id, patient_id, admit_time, discharge_time
  labs.csv           patient_id, encounter_id, time, name, value, unit, ref_low, ref_high,
                     code, vocabulary
  medications.csv    patient_id, encounter_id, time, end_time, name, value, unit
  conditions.csv     patient_id, time, name, code, vocabulary
  notes.csv          patient_id, time, name, text
  recordings.csv     patient_id, time, category, file   (paths relative to this folder)
\end{lstlisting}
The file names choose the categories (\code{labs}, \code{medications}, \code{conditions},
\code{notes}), and the \code{category} column of \file{recordings.csv} chooses each recording's
(\code{icu}, \code{holter}, \code{eeg}, \code{sleep}). The following loop previews and imports
the files:
\begin{lstlisting}[style=shell]
cd study
for t in patients encounters labs medications conditions notes recordings; do
  medorama import-table "$t.csv" --dry-run && medorama import-table "$t.csv" --source study
done
\end{lstlisting}

\paragraph{An SQLite file with the canonical tables.} The simplest database that Medorama
connects to is an SQLite file (or a PostgreSQL schema) whose tables are already the canonical
views. The following SQL creates one:
\begin{lstlisting}[language=SQL]
-- study.sqlite: a database Medorama connects to as it is
CREATE TABLE persons (
  person_key    TEXT PRIMARY KEY,        -- the database's own key
  patient_id    TEXT,                    -- the study's ID (links to imported patients)
  mrn           TEXT,                    -- or the MRN, when the study may hold it
  mrn_authority TEXT DEFAULT '',
  last_name     TEXT DEFAULT '',         -- leave the identifiers empty when not needed
  first_name    TEXT DEFAULT '',
  birth_date    TEXT DEFAULT '',         -- YYYY-MM-DD
  sex           TEXT DEFAULT ''          -- F, M, ...
);
CREATE TABLE visits (
  visit_key  TEXT PRIMARY KEY,
  person_key TEXT NOT NULL REFERENCES persons(person_key),
  start_time TEXT, end_time TEXT,        -- 2024-03-02 07:10, or ISO 8601 with an offset
  unit       TEXT DEFAULT '',
  icu        INTEGER DEFAULT 0
);
CREATE TABLE entries (
  entry_key  TEXT PRIMARY KEY,
  person_key TEXT NOT NULL REFERENCES persons(person_key),
  visit_key  TEXT REFERENCES visits(visit_key),
  category   TEXT NOT NULL,              -- lab, vit, med, dx, proc, note, obs, ...
  time       TEXT NOT NULL,
  end_time   TEXT,
  name       TEXT NOT NULL,
  code       TEXT DEFAULT '', vocabulary TEXT DEFAULT '',
  value_num  REAL, value_text TEXT DEFAULT '', unit TEXT DEFAULT '',
  ref_low    REAL, ref_high REAL, flag TEXT DEFAULT '',
  author     TEXT DEFAULT '', text TEXT DEFAULT ''
);
CREATE INDEX entries_person ON entries(person_key, time);
CREATE INDEX entries_name ON entries(name);
CREATE TABLE loads (loaded_at TEXT);     -- the data's version
\end{lstlisting}
Its views then only select the tables:
\begin{lstlisting}[language=TOML]
# study-views.toml
persons = '''SELECT person_key, patient_id, mrn, mrn_authority, last_name, first_name,
  birth_date, sex FROM persons'''
visits = "SELECT visit_key, person_key, start_time, end_time, unit, icu FROM visits"
entries = '''SELECT entry_key, person_key, visit_key, category, time, end_time, name, code,
  vocabulary, value_num, value_text, unit, ref_low, ref_high, flag, author, text FROM entries'''
version = "SELECT MAX(loaded_at) FROM loads"
\end{lstlisting}
The following commands connect the file, check it, and link its persons:
\begin{lstlisting}[style=shell]
medorama connections add study --kind sql --url sqlite:/data/study/study.sqlite \
    --views study-views.toml
medorama connections check study
medorama connections sync study
\end{lstlisting}
\cmd{connections add} answers with the views it found and the version
(\code{Connected study (sql): views persons, visits, entries; version 2024-03-05}), and
\cmd{sync} with the number of persons linked. From the page, the file must lie inside the data
roots, and it cannot be the catalog's own database.

\paragraph{Choices to make when building one.}
\begin{itemize}
\item \textbf{Identifiers.} A study's own \code{patient\_id} is enough to link its persons to the
  patients of other databases that use the same ID. Hold MRNs, names and dates of birth only when
  they are needed.
\item \textbf{Times.} Choose one time zone for the database and give it as the connection's
  \code{timezone}, or store ISO 8601 times with offsets.
\item \textbf{Codes.} Standard vocabularies (LOINC for results, ICD-10 for diagnoses, RxNorm for
  drugs) let cohorts search by code across databases.
\item \textbf{Categories.} Store the category keys themselves (\code{lab}, \code{med},
  \code{dx}), in lower case.
\item \textbf{Files} (recordings, ECGs, images) cannot be connected. List them in a small table
  (\code{patient\_id}, \code{time}, \code{category}, \code{file}) and import it
  (\cref{chap:importing}).
\item \textbf{Already in OMOP?} Connect it as \code{omop}: there are no views to write.
\end{itemize}

\section{Checking a database}
\label{sec:connecting-checking}

\begin{itemize}
\item \textbf{A connection.} \cmd{medorama connections check NAME} (in the page, \ui{Check})
  asks each view for one row of its canonical columns and the version query for its value, and
  prints how long the database took to answer
  (\code{omop: answers in 41 ms; views persons, visits, entries; version 2024-01-31 / 2023-11-01}).
  A view that does not answer is named, with the database's error. The same check runs when a
  connection is added.
\item \textbf{An import.} \cmd{medorama import-table TABLE -{}-dry-run} reads and checks every row
  without keeping anything (\cref{chap:importing}).
\item \textbf{After either.} The Databases tab lists the database. \ui{Review} lists the rows and
  persons held back and the possible duplicates. A linked patient's page shows each entry in its
  column of the Patient Journey (\cref{chap:patients}); a category that no database provides has
  no column. \cmd{medorama sources list} shows what each source label contributed.
\end{itemize}

\section{The person registry and linking}
\label{sec:connecting-registry}

One person may be in several sources: imported tables, an OMOP instance, a warehouse. The
\emph{registry} is the catalog's patients, each with a random pseudonym, its identifiers, its
merges and its reviewed pairs. A connected database's persons are \emph{linked} to them by the
following rules.
\begin{itemize}
\item \textbf{When.} A person is linked the first time a source names them (in a cohort's search
  or a sync), so that a large database does not have to be copied in.
  \cmd{medorama connections sync NAME} (\ui{Link all}) links every person at once, 2000 at a
  time, and says how many were linked; the others are held for review, or have no valid
  identifier. A sync reads every person in one query, within \code{max\_rows}, so a database with
  more persons than \code{max\_rows} fails to sync. Raise \opt{-{}-max-rows}, or let searches
  link the persons they find.
\item \textbf{How.} A person's identifiers are read as a row of an imported table is, by the same
  rules (\cref{chap:importing}): by patient ID (and its merges), otherwise by MRN of the same
  authority (the view's \code{mrn\_authority}, otherwise the connection's), otherwise as a new
  patient. A person whose patient ID and MRN name two different patients is held for Review, as an
  imported row is. A person with neither is known by the connection's name and its key
  (\code{omop:1234}), as a patient of its own.
\item \textbf{What is copied.} Only what linking needs, kept as the connection's batch: the
  patients it adds, their MRNs and their values (sex, date of birth). Disconnecting removes them,
  as removing an imported database does.
\item \textbf{Merges} move the links with the rest of the patient.
\item \textbf{Never across sources by a guess.} Two sources' persons are the same patient only
  by an identifier that both hold, or by a reviewer's merge. Columns are never mixed between
  sources.
\end{itemize}

\paragraph{On the pages and in cohorts.} A linked patient's page shows the entries of every
source in the same columns. A connected database's visits are listed after the encounters, and
its entries are named by a token (\code{live:<connection>:<token>}, a keyed hash) that holds
neither the entry's key nor its time (\cref{chap:patients}). A criterion is answered by every
source that holds such data, and it can name one source, as it names one database:
\code{omop.lactate > 4} (\cref{chap:cohorts}). A cohort records each source's version when it is
made or refreshed.

\paragraph{Failures.} A source that does not answer leaves a patient's page with what the other
sources gave, and the page names it. A cohort's search over a source that does not answer fails
with that source's name, because a result that silently missed a source would be wrong.

\paragraph{Not there yet.} The following are not supported yet: FHIR servers; SQL Server, Oracle,
Snowflake and BigQuery directly; files held by a connected database (its entries carry no
attachment); pushing value thresholds down to the database when units agree; a connected
database's visits as encounters of their own, and its data counted in the Patient Catalog's
columns; and the notes of a connected database among the notes shown near an entry.

\section{Disconnecting}
\label{sec:connecting-removing}

\cmd{medorama connections remove NAME} (in the page, \ui{Disconnect\dots}) removes the connection,
its links, what linking added to the catalog, and the password file that the page kept. The
database itself is not touched.

\section{Checklist}
\label{sec:connecting-checklist}

\begin{enumerate}
\item The account can only read, and its password is in a file, never in the URL.
\item Every person is named by \code{patient\_id} or \code{mrn}, with \code{mrn\_authority} (in
  the view, or the connection's) when MRNs come from several systems.
\item Every dated row has a time in a readable form, and \code{timezone} is set for local times.
\item Categories are given by their keys, or by names that make them clear.
\item Values are numeric where they are numbers, with their units, and codes come with their
  vocabularies.
\item Identifiers are limited to what linking needs.
\item The underlying tables are indexed on the person, the time, and the name or code.
\item Files are listed in an imported table, inside \code{data\_roots}.
\item A check comes before the real use; \cmd{connections sync} runs for cohorts that must see
  every person; and Review is checked afterwards.
\end{enumerate}
